\section{제안 기법}
\label{sec:제안기법}

제안하는 \method\ 프레임워크는 전역 정보를 통해 오프라인에서 훈련되는 privileged 교사 모델, 로컬 1-hop 정보만을 사용하여 실시간 제어 결정을 학습하는 학생 모델, 그리고 가벼운 nominal 지리적 포워딩과 위험도 예측 기반의 라우팅 분기를 조율하는 예측적 위험 스위칭 모듈로 구성된다.

\subsection{특권적 글로벌 교사 학습}
글로벌 교사 정책 $\pi_T(a \mid s_t; \theta_T)$는 매 라우팅 시점마다 시뮬레이터가 획득한 전역 네트워크 토폴로지 정보 $s_t = (\mathcal{G}_t, p_t, v_t, q_t, d)$ 전체를 바탕으로 최적의 차홉 확률 분포를 모델링한다. 교사의 목적함수는 기대 누적 할인 보상을 극대화하는 것이다.
\begin{equation}
    J(\theta_T) = \mathbb{E}_{\pi_T} \left[ \sum_{t=0}^{T_{\mathrm{ep}}} \gamma^t r_t \right]
\end{equation}
이때 각 홉별 라우팅 보상 $r_t$는 전송 성공 여부, 전송 지연 시간, 에너지 소비 추정치, 홉 수 제약 및 경로 루핑/드롭 실패 페널티를 결합한다.
\begin{equation}
    r_t = \alpha_{\mathrm{succ}}R_{\mathrm{succ}} - \alpha_D D_t - \alpha_E E_t - \alpha_H H_t - \alpha_F F_t
\end{equation}

교사 모델의 매개변수 $\theta_T$는 PPO(Proximal Policy Optimization) 알고리즘을 사용해 최적화한다. 시간 $t$에서의 정책 비율 $\rho_t(\theta_T)$는 다음과 같고,
\begin{equation}
    \rho_t(\theta_T) = \frac{\pi_T(a_t \mid s_t; \theta_T)}{\pi_T(a_t \mid s_t; \theta_T^{\mathrm{old}})}
\end{equation}
가치 네트워크 $V_T(s_t; \phi_T)$를 기반으로 추정한 일반화된 어드밴티지(Advantage) $A_t = \hat{R}_t - V_T(s_t; \phi_T)$를 사용하여 클리핑된 목적함수를 최대화한다.
\begin{equation}
    \mathcal{L}_{\mathrm{PPO}}(\theta_T) = - \mathbb{E}_t \left[ \min \left( \rho_t(\theta_T)A_t, \; \mathrm{clip}(\rho_t(\theta_T), 1-\epsilon, 1+\epsilon)A_t \right) \right]
\end{equation}
가치 함수는 실제 할인된 리턴 $\hat{R}_t$와의 평균 제곱 오차(MSE)를 최소화하도록 훈련된다.
\begin{equation}
    \mathcal{L}_{V}(\phi_T) = \mathbb{E}_t \left[ \left( V_T(s_t; \phi_T) - \hat{R}_t \right)^2 \right]
\end{equation}
최종 교사 네트워크의 총 손실 함수는 다음과 같으며, $\mathcal{H}$는 정책의 탐색을 유도하는 엔트로피 보너스 항이다.
\begin{equation}
    \mathcal{L}_{T}(\theta_T, \phi_T) = \mathcal{L}_{\mathrm{PPO}}(\theta_T) + c_V \mathcal{L}_{V}(\phi_T) - c_H \mathcal{H}(\pi_T(\cdot \mid s_t; \theta_T))
\end{equation}

\subsection{글로벌-로컬 정책 증류}
교사 네트워크의 가중치 $\theta_T^\star$가 수렴하면 이를 동결하고, 로컬 1-hop 정보만을 인풋으로 취하는 학생 네트워크 $\pi_S(a \mid o_u; \theta_S)$를 학습시킨다. 온도 매개변수 $T \ge 1$을 사용해 교사와 학생의 차홉 로짓을 확률 분포로 스케일링한다.
\begin{equation}
    y_T(a) = \mathrm{softmax}\left(\frac{z_T(a)}{T}\right), \quad y_S(a) = \mathrm{softmax}\left(\frac{z_S(a)}{T}\right)
\end{equation}
학생 모델의 학습 손실 함수는 부드러운 분포 모방을 위한 KL Divergence, 교사의 하드 레이블 최적 액션 모방을 위한 Cross Entropy, 그리고 임의의 토폴로지 내에서 유효한 포워딩을 강제하는 최단 경로 오라클 보조 손실 함수를 조합한다.
\begin{equation}
    \mathcal{L}_{\mathrm{KD}}(\theta_S) = \lambda_{\mathrm{KL}} T^2 D_{\mathrm{KL}}(y_T \Vert y_S) + \lambda_{\mathrm{CE}} \mathrm{CE}(a_T, \pi_S) + \lambda_{\mathrm{O}} \mathcal{L}_{\mathrm{oracle}}
\end{equation}

\subsection{예측적 후보 링크 피처 및 P+ 확장}
\method는 무선 채널의 확률적 단절과 급격한 토폴로지 변화 하의 포워딩 위험을 줄이기 위해, 각 후보 이웃 노드 $i \in \mathcal{N}_u(t)$에 대해 기본 물리 피처와 P+ 확장 안정성 피처를 로컬로 평가한다.
\begin{equation}
    x_i^{\mathrm{risk}} = [m_i, \ell_i, q_i, o_i]
\end{equation}
\begin{equation}
    x_i^{+} = [\bar{o}_{i,\mathrm{top}k}, \rho_i, p_{i,\mathrm{keep}}, \eta_i]
\end{equation}
여기서 $m_i$는 수신 RSSI 마진, $\ell_i$는 예측 링크 수명, $q_i$는 버퍼 대기열 여유 용량, $o_i$는 2-hop 방향의 최대 링크 수명이다.
또한, P+ 추가 성분으로 $\bar{o}_{i,\mathrm{top}k}$는 노드 $i$의 outgoing 링크 중 상위 $k$개 링크 수명의 평균값이며, $\rho_i$는 가용한 주변 대체 링크 수(중복성), $p_{i,\mathrm{keep}}$는 채널 변동 모형에 기반한 링크 잔존 확률, 그리고 $\eta_i = 1 - (d_i / R_c)^2$는 상대 거리 $d_i$를 활용한 에너지 효율 대리 지표이다.

\subsection{위험도 점수 및 안전도 평가}
P+ 로컬 피처를 활용하여 각 후보 링크의 위험도 점수(Danger Score) $D_i$와 안전도 효용(Safety Utility) $Q_i$를 공식화한다.
\begin{equation}
    D_i = [g_m - m_i]_+ + [g_\ell - \ell_i]_+ + [g_o - o_i]_+ + [g_k - \bar{o}_{i,\mathrm{top}k}]_+ + 0.5[g_\rho - \rho_i]_+ + [g_p - p_{i,\mathrm{keep}}]_+
\end{equation}
\begin{equation}
    Q_i = m_i + \ell_i + o_i + \bar{o}_{i,\mathrm{top}k} + 0.5\rho_i + p_{i,\mathrm{keep}}
\end{equation}
이때 $g$로 표기된 매개변수들은 각 지표별 안전 임계값(Safety Gate)을 의미하며, $[x]_+ = \max(x, 0)$이다. Nominal 지리 정책에 따른 차홉 행동 $a_N$과 예측 제어 루틴이 추천하는 대안 차홉 행동 $a_P$ 간의 안전도 갭(Safety Gain)은 다음과 같이 정의된다.
\begin{equation}
    G(a_P, a_N) = Q_{a_P} - Q_{a_N}
\end{equation}

\subsection{예측적 위험 스위칭 결정 논리}
기본적으로 가동 효율 극대화를 위해 nominal geographic routing이 출력하는 행동 $a_N$을 우선 사용한다. 그러나 nominal 경로의 안정성이 임계치를 벗어날 때만 예측 스위치 변수 $S \in \{0, 1\}$을 활성화하여 예측 모델 분기 $a_P$를 선택한다.
\begin{equation}
    S = \mathbb{1}\left[ a_N = \dropact \quad \lor \quad D_{a_N} > \tau \quad \lor \quad \left( a_P \neq a_N \; \land \; G(a_P, a_N) > \delta \; \land \; D_{a_P} < D_{a_N} \right) \right]
\end{equation}
여기서 $\tau$는 최대 허용 위험도 임계치이며, $\delta$는 대안 경로 선택 시 보장해야 하는 최소 안전 마진이다. 최종 라우팅 결정 $a^\star$는 다음과 같다.
\begin{equation}
    a^\star = \begin{cases} a_P, & S = 1 \\ a_N, & S = 0 \end{cases}
\end{equation}

\subsection{에너지 보정 및 강제 드롭 억제}
\method는 전송 효율을 물리적으로 뒷받침하기 위해 학생 추론 logit에 두 가지 조정을 추가한다.
첫째, 안전도가 유사한 이웃 노드들이 경합할 경우 송신 거리가 가장 짧고 에너지가 절약되는 노드를 선정할 수 있도록 에너지 보정을 가한다.
\begin{equation}
    z_i^{P+} = z_i^P + \lambda_E \eta_i
\end{equation}
둘째, 로컬 주변에 안전 게이트를 통과하는 전송 노드가 최소 하나라도 존재하는 상황에서는 소극적인 드롭 결정을 내리는 빈도를 제어하기 위해 드롭 억제 페널티를 부여한다.
\begin{equation}
    z_{\dropact}^{P+} = z_{\dropact} - \lambda_D
\end{equation}
이러한 에너지 tie-breaking과 드롭 억제는 네트워크 수준의 소모 전력 및 패킷 생존 성능 향상에 기여한다.
